% TOP_PROBLEM: {"id": 20001071, "title": "Local average outdegree and short directed cycles", "category_id": 2, "category": {"id": 2, "name": "combinatorics", "display_name": "Combinatorics", "description": "Counting problems, graph theory, discrete structures.", "slug": "combinatorics", "order_index": 2, "created_at": "2026-07-31T15:26:25.671Z"}, "status": "partially_solved", "problem_number": "AIM-COMBINATORICS-0196", "set_id": 14, "statusQualification": "States the intended simple loopless version; arbitrary parallel arcs invalidate a degree–girth assertion of this form."}
% FIRST_POSTED: 2026-10-01
% ENTRY_KIND: statement_only
% UnsolvedMath ID 20001071; source catalogue code AIM-COMBINATORICS-0196
% !TeX program = lualatex
% Definition and sources only; this file is not an LLM solution attempt.
\documentclass[11pt,a4paper]{article}
\usepackage[margin=25mm]{geometry}
\usepackage{fontspec}
\setmainfont{lmroman10-regular.otf}[BoldFont=lmroman10-bold.otf,
 ItalicFont=lmroman10-italic.otf,BoldItalicFont=lmroman10-bolditalic.otf]
\usepackage{amsmath,amssymb,mathtools}
\usepackage{unicode-math}
\setmathfont{latinmodern-math.otf}
\AtBeginDocument{\let\setminus\smallsetminus}
\usepackage{xurl,hyperref}
\hypersetup{colorlinks=true,urlcolor=blue}
\setcounter{secnumdepth}{0}
\setlength{\parindent}{0pt}
\setlength{\parskip}{5pt}
\setlength{\emergencystretch}{3em}
\begin{document}
\section*{Local average outdegree and short directed cycles}\label{20001071}
{\small UnsolvedMath ID 20001071 · AIM-COMBINATORICS-0196 · Combinatorics · AIM Workshop Problem Lists}
\subsection{Definitions and mathematical statement}
Let $D=(V,E)$ be a finite simple loopless digraph on $n$ vertices,
with at least one outgoing arc at every vertex. Opposite arcs are
allowed. Write $d^+(v)$ for the number of arcs leaving $v$.
The directed girth $g(D)$ is the length of its shortest directed
cycle; such a cycle exists because following outgoing arcs in a
finite graph must eventually revisit a vertex.
For a positive integer $r$, assume the local degree-sum condition
\[
                       d^+(u)+d^+(v)\ge2r
                              \qquad\text{for every arc }u\to v.
\]
Shen's conjecture asks whether these assumptions imply
\[
                              g(D)\le\left\lceil\frac nr\right\rceil.
\]
The degree sum is imposed on endpoints of every directed edge, not
on all pairs of vertices and not merely on average over the graph.
A vertex may individually have outdegree smaller than $r$, provided
its incident outgoing-edge endpoint sums meet the requirement.
Consequently this is a local average-degree extension of the usual
minimum-outdegree short-cycle question.

The assertion is uniform over $n,r$ and all such digraphs; $r$ is not
the number of vertices or an average degree computed afterwards.
The minimum-outdegree-one assumption is retained independently of the
degree-sum condition, excluding sinks that the latter might not
constrain. Degrees count neighbours in a simple digraph; parallel
arcs would artificially enlarge them without reducing girth.
\subsection{Short English statement}
Does a lower bound on the sum of the endpoint outdegrees of each arc force the same short-cycle bound as a minimum-outdegree condition?
\subsection{Sources}
\begin{itemize}
\item[S1] ULAM.AI and the UnsolvedMath Contributors, supplied problems.json, record 20001071 (AIM-COMBINATORICS-0196), AIM Workshop Problem Lists. Catalogue identity and source transcription; CC BY 4.0.
\url{https://huggingface.co/datasets/ulamai/UnsolvedMath}.
\item[S2] American Institute of Mathematics, The Caccetta-Haggkvist conjecture, item 6. Original workshop question underlying AIM-COMBINATORICS-0196.
\url{https://aimath.org/WWN/caccetta/caccetta.pdf}.
\end{itemize}
\end{document}
